% Options for packages loaded elsewhere
\PassOptionsToPackage{unicode}{hyperref}
\PassOptionsToPackage{hyphens}{url}
\PassOptionsToPackage{space}{xeCJK}
%
\documentclass[
]{article}
\usepackage{amsmath,amssymb}
% \usepackage{lmodern}
\usepackage{iftex}
\ifPDFTeX
  \usepackage[T1]{fontenc}
  \usepackage[utf8]{inputenc}
  \usepackage{textcomp} % provide euro and other symbols
\else % if luatex or xetex
  \usepackage{unicode-math}
  \defaultfontfeatures{Scale=MatchLowercase}
  \defaultfontfeatures[\rmfamily]{Ligatures=TeX,Scale=1}
  \setmainfont[]{DejaVu Serif}
  \setmathfont[]{Asana Math}
  \ifXeTeX
    \usepackage{xeCJK}
    \setCJKmainfont[]{Noto Serif CJK TC}
  \fi
  \ifLuaTeX
    \usepackage[]{luatexja-fontspec}
    \setmainjfont[]{Noto Serif CJK TC}
  \fi
\fi
% Use upquote if available, for straight quotes in verbatim environments
\IfFileExists{upquote.sty}{\usepackage{upquote}}{}
\IfFileExists{microtype.sty}{% use microtype if available
  \usepackage[]{microtype}
  \UseMicrotypeSet[protrusion]{basicmath} % disable protrusion for tt fonts
}{}
\makeatletter
\@ifundefined{KOMAClassName}{% if non-KOMA class
  \IfFileExists{parskip.sty}{%
    \usepackage{parskip}
  }{% else
    \setlength{\parindent}{0pt}
    \setlength{\parskip}{6pt plus 2pt minus 1pt}}
}{% if KOMA class
  \KOMAoptions{parskip=half}}
\makeatother
\usepackage{xcolor}
\IfFileExists{xurl.sty}{\usepackage{xurl}}{} % add URL line breaks if available
\IfFileExists{bookmark.sty}{\usepackage{bookmark}}{\usepackage{hyperref}}
\hypersetup{
  hidelinks,
  pdfcreator={LaTeX via pandoc}}
\urlstyle{same} % disable monospaced font for URLs
\setlength{\emergencystretch}{3em} % prevent overfull lines
\providecommand{\tightlist}{%
  \setlength{\itemsep}{0pt}\setlength{\parskip}{0pt}}
\setcounter{secnumdepth}{-\maxdimen} % remove section numbering
\usepackage[margin=1in]{geometry}
\usepackage{microtype}
\usepackage{hyperref}
\hypersetup{colorlinks=true, linkcolor=blue, urlcolor=blue}

\usepackage{fontspec}
\usepackage{xeCJK}

\setmainfont{DejaVu Serif}
\setCJKmainfont{Noto Serif CJK TC}[BoldFont = Noto Serif CJK TC]

\usepackage{amsmath, amssymb, amsthm}
\ifLuaTeX
  \usepackage{selnolig}  % disable illegal ligatures
\fi

\author{}
\date{}

\begin{document}

我們在 Day11 討論了 MI 協議。但是我們在當時所定義的公鑰是以下程式碼：

\begin{verbatim}
def Public_key(x):

x = T(x)

x = Central_Map_poly(x)

x = S(x)

return x
\end{verbatim}

那這不就代表說，當我們公佈公鑰時，其中的私鑰 T, F, S 都必須一起發出去？

另外，在 Day12 討論了 MI 裡面的 F 是多元二次，所以 P(x) = S(F(T(x)))
是一個多元二次多項式系統。今天我們就使用程式來計算這個公鑰 P(x)
於是我們可以把它安全的發出去。

\hypertarget{ux591aux5143ux4e8cux6b21ux591aux9805ux5f0fux7684ux77e9ux9663ux8868ux793aux6cd5}{%
\section{多元二次多項式的矩陣表示法}\label{ux591aux5143ux4e8cux6b21ux591aux9805ux5f0fux7684ux77e9ux9663ux8868ux793aux6cd5}}

首先根據 Day12 的內容，我們知道公鑰 P 可以寫成：

\$\$

\mathcal{P} =

\left\{

\begin{aligned}

& p^{(1)}\left(x_1, x_2, \ldots, x_n\right)= \sum_{i=1}^n \sum_{j=i}^n p_{i j}^{(1)} x_i x_j+\sum_{i=1}^n p_i^{(1)} x_i+p_0^{(1)} \\

& p^{(2)}\left(x_1, x_2, \ldots, x_n\right)= \sum_{i=1}^n \sum_{j=i}^n p_{i j}^{(2)} x_i x_j+\sum_{i=1}^n p_i^{(2)} x_i+p_0^{(2)} \\

& \vdots \\

& p^{(m)}\left(x_1, x_2, \ldots, x_n\right)=\sum_{i=1}^n \sum_{j=i}^n p_{i j}^{(m)} x_i x_j+\sum_{i=1}^n p_i^{(m)} x_i+p_0^{(m)}

\end{aligned}

\right.

\$\$

其中係數都是固定的。

如果想要安全地把私鑰送出去，那我們就得乖乖的把上面這些係數通通算出來，然後發送這些係數出去。

為了讓程式好寫，我們在這裡使用矩陣表示法：

首先先看其中第 k 個多元二次多項式：

\$\$

p\^{}\{(k)\}\left(x\_1, x\_2, \ldots, x\_n\right)=

\sum\emph{\{i=1\}\^{}n \sum}\{j=i\}\^{}n p\_\{i j\}\^{}\{(k)\} x\_i
x\_j+

\sum\_\{i=1\}\^{}n p\_i\^{}\{(k)\} x\_i+

p\_0\^{}\{(k)\}

\$\$

如果把 x 寫成行向量

\$\$

x =

\begin{bmatrix}

x_1 \\ \vdots \\ x_n

\end{bmatrix}

\$\$

則有

\$\$

\begin{aligned}

\sum_{i=1}^n \sum_{j=i}^n p_{i j} x_i x_j

&=

\begin{bmatrix}

x_1 & \cdots & x_n

\end{bmatrix}

\begin{bmatrix}

p_{11} & p_{12} & \cdots & p_{1n} \\

0 & p_{22} & \cdots & p_{2n} \\

\vdots & \vdots & \ddots & \vdots \\

0 & 0 & \cdots & p_{nn} \\

\end{bmatrix}

\begin{bmatrix}

x_1 \\ \vdots \\ x_n

\end{bmatrix}

\\

&=x^{T}Qx

\end{aligned}

\$\$

\$\$

\begin{aligned}

\sum_{i=1}^n p_i x_i &=

\begin{bmatrix}

p_{1} & p_{2} & \cdots & p_{n}

\end{bmatrix}

\begin{bmatrix}

x_1 \\ \vdots \\ x_n

\end{bmatrix}

\\

&= L x

\end{aligned}

\$\$

於是我們可以寫出

\$\$

p\^{}\{(k)\}\left(x\_1, x\_2, \ldots, x\_n\right)=

x\textsuperscript{\{T\}Q}\{(k)\}x+L\^{}\{(k)\}x +C\^{}\{(k)\}

\$\$

其中

\$\$

Q\^{}\{(k)\} =

\begin{bmatrix}

p_{11}^{(k)} & p_{12}^{(k)} & \cdots & p_{1n} ^{(k)}\\

0 & p_{22}^{(k)} & \cdots & p_{2n}^{(k)} \\

\vdots & \vdots & \ddots & \vdots \\

0 & 0 & \cdots & p_{nn}^{(k)} \\

\end{bmatrix}

\quad

L\^{}\{(k)\} =

\begin{bmatrix}

p_{1} ^{(k)}& p_{2} ^{(k)}& \cdots & p_{n} ^{(k)}

\end{bmatrix}

,

\quad C\^{}\{(k)\} = p\_\{0\}\^{}\{(k)\}

\$\$

結論：

我們要計算出每個矩陣 Q, L, 以及常數項 C ，然後就可以當作公鑰發出去。

\hypertarget{mi-ux7cfbux7d71ux88e1ux7684ux7279ux6b8aux89c0ux5bdf}{%
\subsection{MI
系統裡的特殊觀察}\label{mi-ux7cfbux7d71ux88e1ux7684ux7279ux6b8aux89c0ux5bdf}}

因為我們在 MI 系統內要設定 q = 2，也就是係數是 Z\_2 裡面的元素，因此

\$\$

x\_i\^{}2 = x\_i

\$\$

所以所有的線性項都可以當作二次項

\$\$

\sum\_\{i=1\}\^{}n p\_i x\_i =

\sum\_\{i=1\}\^{}n p\_i x\_i\^{}2

\$\$

也因此，對 MI 系統來說，我們可以把 p\^{}\{(k)\} 寫成

\$\$

p\^{}\{(k)\}(x\_1,\ldots,x\_n) = x\^{}T Q\^{}\{(k)\} x +C\^{}\{(k)\}

\$\$

\hypertarget{ux8a08ux7b97-mi-ux7cfbux7d71ux7684ux516cux9470ux4fc2ux6578}{%
\section{計算 MI
系統的公鑰係數}\label{ux8a08ux7b97-mi-ux7cfbux7d71ux7684ux516cux9470ux4fc2ux6578}}

接下來，我們使用 SageMath 來實作這個矩陣表示法，並計算公鑰。

\hypertarget{ux6e96ux5099ux5de5ux4f5cux518dux505aux4e00ux6b21-mi}{%
\subsection{準備工作（再做一次 MI
）}\label{ux6e96ux5099ux5de5ux4f5cux518dux505aux4e00ux6b21-mi}}

首先，讓我們定義參數和隨機仿射映射

\begin{verbatim}
q = 2

n = 5 # degree of g(x)

  

R = quotient(ZZ, q*ZZ)

R_poly = PolynomialRing(R, 'x')

  

g = x^5 + x^3 + 1

R_poly_quotient = quotient(R_poly, g)

  

theta = 2

h = (xgcd((q^n)-1, (q^theta)+1))[2]

h = h + (q^n-1)

  

# 生成隨機仿射映射

def RandomAffineMapGenerator(n, R):

while True:

A = random_matrix(R, n, n)

if A.is_invertible():

break

b = random_vector(R, n)

  

def AffineMap(v):

v = vector(v)

v = A * v + b

return v.list()

  

def InverseMap(v):

v = vector(v)

v = A.inverse() * (v - b)

return v.list()

  

return AffineMap, InverseMap, A, b

  

S, S_inv, A_S, b_S = RandomAffineMapGenerator(n, R)

T, T_inv, A_T, b_T = RandomAffineMapGenerator(n, R)
\end{verbatim}

\begin{verbatim}
def Central_Map_poly(X):

X = R_poly_quotient(X)

return (X^(q^theta + 1)).list()

  

def Central_Map_poly_inv(X):

X = R_poly_quotient(X)

return (X^h).list()

  

def Public_key(x):

x = T(x)

x = Central_Map_poly(x)

x = S(x)

return x
\end{verbatim}

\hypertarget{ux958bux59cbux8a08ux7b97ux4fc2ux6578}{%
\subsection{開始計算係數}\label{ux958bux59cbux8a08ux7b97ux4fc2ux6578}}

我們先把每個 p\^{}\{(k)\} 的常數項給算出來，並記錄在 C 這個 list 裡面

\begin{verbatim}
zero_vector = [0] * n

C = Public_key(zero_vector)
\end{verbatim}

另外定義一個扣掉常數項之後的公鑰映射

\begin{verbatim}
def Homogeneous_Public_key(x):

x = Public_key(x)

for i in range(n):

x[i] = x[i] - C[i]

return x
\end{verbatim}

定義這個之後我們就有

\$\$

\text{Homogeneous\_Public\_key}(x) = {[}x\^{}T Q\^{}\{(1)\} x, x\^{}T
Q\^{}\{(2)\} x, \ldots, x\^{}T Q\^{}\{(n)\} x{]}

\$\$

也就是

\$\$

\text{Homogeneous\_Public\_key}(x) {[}k{]} = x\^{}T Q\^{}\{(k)\} x

\$\$

於是就可以很方便計算 Q\^{}\{(k)\}.

我們在程式碼裡面使用 Q 來儲存所有的 Q\^{}\{(k)\}：

\$\$

Q = {[} Q\^{}\{(1)\}, Q\^{}\{(2)\},\ldots, Q\^{}\{(n)\}{]}

\$\$

\begin{verbatim}
# 先初始化為零，把形狀做好

Q = [[[0 for _ in range(n)] for _ in range(n)] for _ in range(n)]
\end{verbatim}

先計算對角線項：

原理是因為如果 e\_i = (0,\ldots,1,\ldots0) （只有在第 i 的位置是 1
其他都是零）那

\$\$

e\_i\^{}T Q\^{}\{(k)\}e\_i = Q\^{}\{(k)\}\_\{ii\}

\$\$

\begin{verbatim}
# 計算對角線項

for i in range(n):

for k in range(n):

e_i = [0 for _ in range(n)]

e_i[i] = 1

Q[k][i][i] = Homogeneous_Public_key(e_i)[k]
\end{verbatim}

接著計算非對角線項：

原理是因為如果 e\_ij = (0,\ldots,1,\ldots,1,\ldots0) （只有在第 i, j
的位置是 1 其他都是零）

則

\$\$

e\_\{ij\}\^{}T Q\^{}\{(k)\}e\_\{ij\} =
Q\^{}\{(k)\}\emph{\{ii\}+Q\textsuperscript{\{(k)\}\emph{\{ij\}+Q\^{}\{(k)\}}\{ji\}+Q}\{(k)\}}\{jj\}

\$\$

但因為我們所設定的 Q 是上三角矩陣，所以 Q\_ji = 0

因此我們知道

\$\$

\begin{aligned}

Q^{(k)}_{ij} &= e_{ij}^T Q^{(k)}e_{ij} - Q^{(k)}_{ii} - Q^{(k)}_{jj}\\

&= \text{Homogeneous\_Public\_key}(e_{ij}) [k] - Q[k][i][i] - Q[k][j][j]

\end{aligned}

\$\$

\begin{verbatim}
# 計算非對角線項

for k in range(n):

for i in range(n):

for j in range(i+1,n):

e_ij = [0 for _ in range(n)]

e_ij[i] = 1

e_ij[j] = 1

  

Q[k][i][j] = Homogeneous_Public_key(e_ij)[k] - (Q[k][i][i]) - Q[k][j][j]

  
\end{verbatim}

好！幾乎完成了！現在已經可以把 Q = {[}Q\^{}\{(1)\}, \ldots,
Q\^{}\{(n)\}{]} 以及常數項 C 釋出。

\hypertarget{ux9a57ux8b49ux6b63ux78baux6027}{%
\subsection{驗證正確性}\label{ux9a57ux8b49ux6b63ux78baux6027}}

為了驗證正確性，我們進行以下的驗證。

首先把 Q 裡面的每個矩陣 Q\^{}\{(k)\} 都轉化為 SageMath
裡面的矩陣類別，後面的乘法就可以寫得很簡單：

\begin{verbatim}
for k in range(n):

Q[k] = matrix(Q[k])
\end{verbatim}

於是可以做出以下根據矩陣 Q\^{}\{(k)\} 以及常數項 C 所做的公鑰：

\begin{verbatim}
def Publishable_Public_Key(x):

x = (Matrix([[x[i]] for i in range(n)]))

result = []

for k in range(n):

result.append(x.T * Q[k] * x + C[k])

  
  

result = [result[i][0][0] for i in range(n)]
\end{verbatim}

並進行驗證

\begin{verbatim}
for _ in range(100):

x = [R(randint(0,1)) for i in range(5)]

result = Public_key(x)

result2 = Publishable_Public_Key(x)

if not (result == result2):

print("WRONG Public Key")

break
\end{verbatim}

\end{document}
